\sectionkicker{Source-backed technical notes}
\subsection*{Frontier Models \& Voice — モデル更新が利用面へ広がった5月: Technical Notes}
本欄は記事本文で圧縮した一次資料上の情報を、比較・再検証しやすい形で整理したものである。時系列、確認済みの主張、留意点、一次資料URLを掲載する。
\medskip
\subsection*{Theme at a glance}
\addcontentsline{toc}{subsection}{Theme at a glance}
\begin{center}
\footnotesize
\begin{tabularx}{\linewidth}{@{}p{0.20\linewidth}p{0.10\linewidth}p{0.14\linewidth}X@{}}
\toprule
資料 & 位置づけ & 種別 & 時系列 \\
\midrule
GPT-5.5 Instant: smarter, clearer, and more personalized & 主要資料 & モデル更新 & 2026-05-05 (OFFICIAL\_PUBLICATION) \\
Anthropic May 2026 technical releases & 主要資料 & OTHER & 2026-05-28 (モデル\_RELEASE) \\
Gemini API May 2026 lifecycle & 主要資料 & OTHER & 2026-05-19 (モデル\_RELEASE); 2026-05-19 (AGENT\_RELEASE); 2026-05-07 (モデル\_RELEASE); 2026-05-05 (API\_UPDATE); 2026-05-28 (モデル\_RELEASE) \\
Google DeepMind May 2026 releases & 主要資料 & OTHER & 2026-05 (モデル\_RELEASE); 2026-05 (モデル\_RELEASE); 2026-05 (研究\_RELEASE); 2026-05 (研究\_RELEASE) \\
Advancing voice intelligence with new models in the API & 主要資料 & API & 2026-05-07 (OFFICIAL\_PUBLICATION) \\
Alibaba Model Studio May 2026 model lifecycle & 補足資料 & OTHER & 2026-05-21 (モデル\_RELEASE); 2026-05-25 (モデル\_RELEASE); 2026-05-19 (モデル\_RELEASE) \\
How OpenAI delivers low-latency voice AI at scale & 補足資料 & OTHER & 2026-05-04 (OFFICIAL\_PUBLICATION) \\
\bottomrule
\end{tabularx}
\end{center}
\normalsize

\begin{technicalnote}{GPT-5.5 Instant: smarter, clearer, and more personalized}{主要資料}
\begin{tabularx}{\linewidth}{@{}>{\bfseries}p{0.22\linewidth}X@{}}
組織 & OpenAI \\
種別 & モデル更新 \\
時系列 & 2026-05-05 (OFFICIAL\_PUBLICATION) \\
\end{tabularx}
\smallskip
{\bfseries 一次資料から整理したtechnical points}
\begin{itemize}[leftmargin=1.5em,itemsep=0.35em]
\item \textbf{一次情報で確認できる事実}: The preserved primary source identifies “GPT-5.5 Instant: smarter, clearer, and more personalized” on 2026-05-05.
\item \textbf{Vendor claim}: OpenAI's official feed records GPT-5.5 Instant as a May 5 update to ChatGPT's default model and describes improvements in answer quality, hallucination reduction, and personalization controls.
\end{itemize}
{\bfseries 読む際の境界}
\begin{itemize}[leftmargin=1.5em,itemsep=0.35em]
\item \textbf{分析上の留意点}: The preserved first-party feed item is authoritative for publication chronology and vendor-described scope; detailed capability, benchmark, or safety claims were not independently reproduced.
\end{itemize}
{\bfseries 一次資料}
\begingroup\sloppy
\begin{itemize}[leftmargin=1.5em,itemsep=0.25em]
\item {\scriptsize\url{https://openai.com/index/gpt-5-5-instant}}
\end{itemize}
\endgroup
\end{technicalnote}
\medskip

\begin{technicalnote}{Anthropic May 2026 technical releases}{主要資料}
\begin{tabularx}{\linewidth}{@{}>{\bfseries}p{0.22\linewidth}X@{}}
組織 & Anthropic \\
種別 & OTHER \\
時系列 & 2026-05-28 (モデル\_RELEASE) \\
\end{tabularx}
\smallskip
{\bfseries 一次資料から整理したtechnical points}
\begin{itemize}[leftmargin=1.5em,itemsep=0.35em]
\item \textbf{一次情報で確認できる事実}: The preserved primary source identifies “anthropic-news”.
\item \textbf{Vendor claim}: The Anthropic newsroom snapshot records the May 28 introduction of Claude Opus 4.8, described by Anthropic as an Opus-class upgrade for coding, agentic tasks, professional work, and long-running work.
\end{itemize}
{\bfseries 読む際の境界}
\begin{itemize}[leftmargin=1.5em,itemsep=0.35em]
\item \textbf{分析上の留意点}: The preserved first-party index is authoritative for the listed chronology, but capability/benchmark descriptions remain vendor claims and month-only dates retain month precision.
\end{itemize}
{\bfseries 一次資料}
\begingroup\sloppy
\begin{itemize}[leftmargin=1.5em,itemsep=0.25em]
\item {\scriptsize\url{https://www.anthropic.com/news}}
\end{itemize}
\endgroup
\end{technicalnote}
\medskip

\begin{technicalnote}{Gemini API May 2026 lifecycle}{主要資料}
\begin{tabularx}{\linewidth}{@{}>{\bfseries}p{0.22\linewidth}X@{}}
組織 & Google \\
種別 & OTHER \\
時系列 & 2026-05-19 (モデル\_RELEASE); 2026-05-19 (AGENT\_RELEASE); 2026-05-07 (モデル\_RELEASE); 2026-05-05 (API\_UPDATE); 2026-05-28 (モデル\_RELEASE) \\
\end{tabularx}
\smallskip
{\bfseries 一次資料から整理したtechnical points}
\begin{itemize}[leftmargin=1.5em,itemsep=0.35em]
\item \textbf{一次情報で確認できる事実}: The preserved primary source identifies “google-gemini-api-release-notes”.
\item \textbf{Vendor claim}: The Gemini API changelog records May releases including Gemini 3.5 Flash GA and Managed Agents public preview on May 19, Gemini 3.1 Flash-Lite GA on May 7, multimodal File Search on May 5, and native visual model GA releases on May 28.
\end{itemize}
{\bfseries 読む際の境界}
\begin{itemize}[leftmargin=1.5em,itemsep=0.35em]
\item \textbf{分析上の留意点}: The preserved first-party index is authoritative for the listed chronology, but capability/benchmark descriptions remain vendor claims and month-only dates retain month precision.
\end{itemize}
{\bfseries 一次資料}
\begingroup\sloppy
\begin{itemize}[leftmargin=1.5em,itemsep=0.25em]
\item {\scriptsize\url{https://ai.google.dev/gemini-api/docs/changelog}}
\end{itemize}
\endgroup
\end{technicalnote}
\medskip

\begin{technicalnote}{Google DeepMind May 2026 releases}{主要資料}
\begin{tabularx}{\linewidth}{@{}>{\bfseries}p{0.22\linewidth}X@{}}
組織 & Google DeepMind \\
種別 & OTHER \\
時系列 & 2026-05 (モデル\_RELEASE); 2026-05 (モデル\_RELEASE); 2026-05 (研究\_RELEASE); 2026-05 (研究\_RELEASE) \\
\end{tabularx}
\smallskip
{\bfseries 一次資料から整理したtechnical points}
\begin{itemize}[leftmargin=1.5em,itemsep=0.35em]
\item \textbf{一次情報で確認できる事実}: The preserved primary source identifies “google-deepmind-blog”.
\item \textbf{Vendor claim}: The DeepMind blog snapshot lists May 2026 entries for Gemini 3.5, Gemini Omni, Gemini for Science, and Co-Scientist, providing a month-level primary chronology across frontier models and AI-for-science.
\end{itemize}
{\bfseries 読む際の境界}
\begin{itemize}[leftmargin=1.5em,itemsep=0.35em]
\item \textbf{分析上の留意点}: The preserved first-party index is authoritative for the listed chronology, but capability/benchmark descriptions remain vendor claims and month-only dates retain month precision.
\end{itemize}
{\bfseries 一次資料}
\begingroup\sloppy
\begin{itemize}[leftmargin=1.5em,itemsep=0.25em]
\item {\scriptsize\url{https://deepmind.google/blog/}}
\end{itemize}
\endgroup
\end{technicalnote}
\medskip

\begin{technicalnote}{Advancing voice intelligence with new models in the API}{主要資料}
\begin{tabularx}{\linewidth}{@{}>{\bfseries}p{0.22\linewidth}X@{}}
組織 & OpenAI \\
種別 & API \\
時系列 & 2026-05-07 (OFFICIAL\_PUBLICATION) \\
\end{tabularx}
\smallskip
{\bfseries 一次資料から整理したtechnical points}
\begin{itemize}[leftmargin=1.5em,itemsep=0.35em]
\item \textbf{一次情報で確認できる事実}: The preserved primary source identifies “Advancing voice intelligence with new models in the API” on 2026-05-07.
\item \textbf{Vendor claim}: OpenAI's official feed records new realtime voice models in the API for speech reasoning, translation, and transcription.
\end{itemize}
{\bfseries 読む際の境界}
\begin{itemize}[leftmargin=1.5em,itemsep=0.35em]
\item \textbf{分析上の留意点}: The preserved first-party feed item is authoritative for publication chronology and vendor-described scope; detailed capability, benchmark, or safety claims were not independently reproduced.
\end{itemize}
{\bfseries 一次資料}
\begingroup\sloppy
\begin{itemize}[leftmargin=1.5em,itemsep=0.25em]
\item {\scriptsize\url{https://openai.com/index/advancing-voice-intelligence-with-new-models-in-the-api}}
\end{itemize}
\endgroup
\end{technicalnote}
\medskip

\begin{technicalnote}{Alibaba Model Studio May 2026 model lifecycle}{補足資料}
\begin{tabularx}{\linewidth}{@{}>{\bfseries}p{0.22\linewidth}X@{}}
組織 & Alibaba Cloud \\
種別 & OTHER \\
時系列 & 2026-05-21 (モデル\_RELEASE); 2026-05-25 (モデル\_RELEASE); 2026-05-19 (モデル\_RELEASE) \\
\end{tabularx}
\smallskip
{\bfseries 一次資料から整理したtechnical points}
\begin{itemize}[leftmargin=1.5em,itemsep=0.35em]
\item \textbf{一次情報で確認できる事実}: The preserved primary source identifies “alibaba-model-studio-model-lifecycle”.
\item \textbf{Vendor claim}: The official Model Studio lifecycle snapshot contains multiple May 2026 model events, including qwen3.7-max on May 21, qwen3.7-max-preview on May 25, and qwen3.5-livetranslate-flash-realtime on May 19.
\end{itemize}
{\bfseries 読む際の境界}
\begin{itemize}[leftmargin=1.5em,itemsep=0.35em]
\item \textbf{分析上の留意点}: The preserved first-party index is authoritative for the listed chronology, but capability/benchmark descriptions remain vendor claims and month-only dates retain month precision.
\end{itemize}
{\bfseries 一次資料}
\begingroup\sloppy
\begin{itemize}[leftmargin=1.5em,itemsep=0.25em]
\item {\scriptsize\url{https://www.alibabacloud.com/help/en/model-studio/newly-released-models}}
\end{itemize}
\endgroup
\end{technicalnote}
\medskip

\begin{technicalnote}{How OpenAI delivers low-latency voice AI at scale}{補足資料}
\begin{tabularx}{\linewidth}{@{}>{\bfseries}p{0.22\linewidth}X@{}}
組織 & OpenAI \\
種別 & OTHER \\
時系列 & 2026-05-04 (OFFICIAL\_PUBLICATION) \\
\end{tabularx}
\smallskip
{\bfseries 一次資料から整理したtechnical points}
\begin{itemize}[leftmargin=1.5em,itemsep=0.35em]
\item \textbf{一次情報で確認できる事実}: The preserved primary source identifies “How OpenAI delivers low-latency voice AI at scale” on 2026-05-04.
\item \textbf{Vendor claim}: OpenAI's official feed describes a rebuilt WebRTC stack for low-latency, globally scaled realtime voice AI and conversational turn-taking.
\end{itemize}
{\bfseries 読む際の境界}
\begin{itemize}[leftmargin=1.5em,itemsep=0.35em]
\item \textbf{分析上の留意点}: The preserved first-party feed item is authoritative for publication chronology and vendor-described scope; detailed capability, benchmark, or safety claims were not independently reproduced.
\end{itemize}
{\bfseries 一次資料}
\begingroup\sloppy
\begin{itemize}[leftmargin=1.5em,itemsep=0.25em]
\item {\scriptsize\url{https://openai.com/index/delivering-low-latency-voice-ai-at-scale}}
\end{itemize}
\endgroup
\end{technicalnote}
\medskip
